\documentclass[12pt]{article}
\usepackage{blindtext}
\usepackage[T1]{fontenc}
\usepackage[utf8]{inputenc}
\usepackage{amsmath}
\usepackage{amsfonts} 
\usepackage{color}
\usepackage{graphicx}
\usepackage{vmargin}
\usepackage{amssymb}
\usepackage[spanish]{babel}


\renewcommand{\abstractname}{Descripción}   % Cambiamos el nombre del abstract
\renewcommand{\contentsname}{Contenidos}   % Cambiamos el nombre de contents


\setlength{\jot}{12pt}

\setmargins{2.5cm}       % margen izquierdo
{1.5cm}                        % margen superior
{16.5cm}                      % anchura del texto
{23.42cm}                    % altura del texto
{10pt}                           % altura de los encabezados
{1cm}                           % espacio entre el texto y los encabezados
{0pt}                             % altura del pie de página
{2cm}                           % espacio entre el texto y el pie de página


\renewcommand{\abstractname}{Descripción}   % Cambiamos el nombre del abstract

\makeatletter
\newcommand*{\bigcdot}{}% Check if undefined
\DeclareRobustCommand*{\bigcdot}{%
  \mathbin{\mathpalette\bigcdot@{}}%
}
\newcommand*{\bigcdot@scalefactor}{.5}
\newcommand*{\bigcdot@widthfactor}{1.15}
\newcommand*{\bigcdot@}[2]{%
  % #1: math style
  % #2: unused
  \sbox0{$#1\vcenter{}$}% math axis
  \sbox2{$#1\cdot\m@th$}%
  \hbox to \bigcdot@widthfactor\wd2{%
    \hfil
    \raise\ht0\hbox{%
      \scalebox{\bigcdot@scalefactor}{%
        \lower\ht0\hbox{$#1\bullet\m@th$}%
      }%
    }%
    \hfil
  }%
}
\makeatother

\title{Tarea 1}
\author{Jesus Eduardo Loera Casas 1898887}
\date{\today}

%%%%%%%%%%%%%%%%%%%%%%%%%%%%%%%%%%%%%%%%%%%%%%%%%%%%%%%%%%%%%%%%%%%%%%%%%%%%
%%%%%%%%%%%%%%%%%%%%%%%%%%%%%%%%%%%%%%%%%%%%%%%%%%%%%%%%%%%%%%%%%%%%%%%%%%%%

\begin{document}

\begin{titlepage}
    \centering
    {\includegraphics[width=0.2\textwidth]{FCFM.png}\par}
    \vspace{1cm}
    {\bfseries\LARGE Universidad Autonoma de Nuevo León \par}
    \vspace{1cm}
    {\scshape\Large Facultad de Ciencias Físico Matemáticas \par}
    \vspace{3cm}
    {\scshape\Huge Tarea 2  \par}
    \vspace{3cm}
    {\itshape\Large Ecuaciones de movimiento \par}
    \vfill
    {\Large Autores: \par}
    \vfill
    {\Large Jesús Eduardo Loera Casas 1898887 \par}
    \vfill{\Large Cesar Efrén Valladares Rocha 1841555 \par}
    \vfill
    {\Large Vrani Chavez Islas 1990044 \par}
    \vfill
    {\Large \today \par}
\end{titlepage}

\tableofcontents			% Índice

\begin{center}
	\rule[0mm]{150mm}{0.1mm}		% Para dibujar una linea horizontal de
									% [elevación]{longitud}{grosor}
	\end{center}
	
	
\begin{abstract}		% ABSTRACT
  
\noindent 				% Anula la sangria de este parrafo
En este documento nuestro equipo presenta el Proyecto 2 del curso de mecánica teórica, 
donde encontramos la ecuación de Euler-Lagrange para una variable independiente y dos dependientes
con ligadura.
\end{abstract}
	
\begin{center}
	\rule[0mm]{150mm}{0.1mm}
	\end{center}

\section{Problema del oscilador armónico}	
    
Encontrar las ecuaciones de movimiento y velocidad respecto al tiempo 
del siguiente oscilador armónico:

\vspace*{0.3 cm}

\begin{center}
    \begin{circuitikz}[american voltages]
        \draw (0,0) to [C=$C$] (4,0) -- (4,2) to [L, l=$L$] (0,2) -- (0,0);
    \end{circuitikz}
\end{center}

Por la ley de Kirchhoff:

\begin{align*}
    L \frac{dI}{dt} + \frac{1}{C} \int Idt = 0
\end{align*}

Donde $ I(t)=\dot{q} $        
                            
\section{Solución}
    \input{solución.tex}

\section{Ecuaciones de movimiento}
    
Por lo tanto la ecuación del movimiento del oscilador armónico simple es:

\begin{itemize}
    \item $x(t) = C_{1}Cos(\omega_{0}t) + C_{2}Sen(\omega_{0}t)$
\end{itemize}

Que puede reescribirse como:

\begin{itemize}
    \item $x(t) = ACos(\omega_{0}t - \phi)$
    \item $x(t) = ASen(\omega_{0}t - \delta)$
\end{itemize}

\end{document}